% Guia do professor — Capítulo 14: Fundamentos de Tempestades
\section{Capítulo 14 --- Fundamentos de Tempestades}

\subsection*{Visão geral e ritmo}

O capítulo em que a termodinâmica dos Caps.~3--5 vira drama: CAPE, CIN,
gatilhos e o papel organizador do cisalhamento. A narrativa dos
``quatro ingredientes'' segura as duas aulas; o clímax é entender por
que o hodógrafo — e não o CAPE — decide o tipo de tempestade.
\textbf{Ritmo: 2 aulas de 2\,h.}

\begin{itemize}
  \item \textbf{Aula 1} — receita de 4 ingredientes; Skew-$T$ com CAPE
        e CIN (Caso~1 ao vivo); $w = \sqrt{2CAPE}$ e o fator
        $\sim$1/2 real; a tarde que estoura a tampa (N1).
  \item \textbf{Aula 2} — hodógrafos e os três tipos (Caso~3);
        supercélula em corte; piscina fria, frente de rajada e RKW
        (Casos~4); ciclo de vida e MCCs sul-americanos.
\end{itemize}

\subsection*{Objetivos de aprendizagem}

\begin{enumerate}
  \item calcular CAPE/CIN/LCL/LFC/EL de uma sondagem (MetPy e à mão no
        Skew-$T$);
  \item converter CAPE em $w_{max}$ e justificar o fator $\sim$1/2
        observado;
  \item classificar o regime convectivo pelo cisalhamento 0--6\,km e
        pelo BRN;
  \item explicar a auto-propagação via piscina fria (corrente de
        densidade) e o equilíbrio RKW;
  \item reconhecer as partes da supercélula (updraft, FFD/RFD,
        bigorna, overshooting top).
\end{enumerate}

\subsection*{Pré-requisitos a reativar}

Skew-$T$ e método da parcela (Cap.~5, essencial!); $T_v$ com carga de
água (Cap.~1); bulbo úmido e evaporação (Cap.~4); Bowen e o
combustível da tarde (Cap.~3); gatilhos frontais (Cap.~12).

\subsection*{Armadilhas conceituais frequentes}

\begin{itemize}
  \item \textbf{``CAPE grande $=$ tempestade severa''}: sem
        cisalhamento, CAPE de 3000 vira um pulso de 30\,min. Martelar
        o papel do hodógrafo.
  \item \textbf{CIN como vilão}: o CIN moderado é o que PERMITE a
        explosão da tarde — sem tampa não há panela de pressão.
  \item \textbf{Sentidos no HS}: hodógrafo útil curva ANTI-horário com
        a altura; mesociclones giram horário. Espelhar toda figura
        importada do HN.
  \item \textbf{$w = \sqrt{2CAPE}$ como previsão}: é teto teórico;
        entranhamento e água cobram metade (N2 mostra).
  \item \textbf{Confundir frente de rajada com frente fria}: mesma
        física de corrente de densidade (Cap.~12!), escalas e vidas
        diferentes.
\end{itemize}

\subsection*{Demonstração de baixo custo}

Corrente de densidade caseira: água salgada tingida despejada num
aquário de água doce — a ``frente de rajada'' com o nariz elevado
aparece perfeita. Vídeos de timelapse de supercélula (rotação visível
do updraft) valem 10 minutos de aula.

\subsection*{Problemas sugeridos do livro (exercícios do Cap.~14)}

Aquecimento: CAPE/CIN de sondagens dadas; $w_{max}$ e comparação com
updrafts típicos. Aprofundamento: BRN e classificação; velocidade de
frente de rajada. Síntese: ``por que o fim de tarde de verão é a hora
das tempestades no interior do Brasil?''.
\emph{Confira a numeração na sua cópia (v1.02b).}

\subsection*{Uso do notebook \texttt{cap14\_tempestades.ipynb}}

\begin{center}
\begin{tabular}{lll}
\hline
Caso & Conteúdo & Momento sugerido \\
\hline
1 & CAPE/CIN com MetPy (Skew-$T$ sombreado) & Aula 1, após parcela \\
2 & $w_{max}(\mathrm{CAPE})$ e o fator 1/2 & Aula 1 \\
3 & hodógrafos: única × multi × super & Aula 2, abertura \\
4 & piscina fria: $c(\Delta\theta_v, H)$ + downdraft & Aula 2 \\
\hline
\end{tabular}
\end{center}

O Caso~1 é o gancho perfeito para rodar ao vivo com a sondagem do dia
(Wyoming, SBMT 12\,UTC) e decidir com a turma ``vai ter tempestade
hoje?''.

\subsection*{Exercícios complementares e soluções}

Nas notas: T1--T5 e N1--N4. Soluções em \texttt{cap14\_solucoes.ipynb}
(\emph{não distribuir}): T2 com \texttt{sympy} (CAPE do triângulo);
N3 implementa SRH com movimento de tempestade do HS. Pares:
T1$\leftrightarrow$N2, T4$\leftrightarrow$N4, T5$\leftrightarrow$N3.

\subsection*{Conexões com o resto do curso}

Método da parcela e estabilidade (Cap.~5); microfísica da chuva e do
granizo (Cap.~7); assinaturas de radar (Cap.~8); gatilhos frontais e
linha seca (Cap.~12); os perigos (Cap.~15); convecção tropical e
furacões (Cap.~16); parametrização de convecção nos modelos (Cap.~20).
